% Vista científica exacta materializada desde una rama condicional.
% Fuente: source/demostraciones_primarias/09_06h_atlas_familias_rev6.tex; rama: HMTIncludeAtlasStructural.
% No modifica el contenido matemático de la rama de origen.
\chapter[Clasificación nonádica de familias]{Clasificación nonádica de familias de trayectoria}
\label{chap:atlas-genealogia-masas}
\index{atlas nonádico}
\index{familia de ruta}
\index{multisección}
\index{sector nulo}

El atlas no comienza con una lista de partículas ni con una tabla de masas.
Comienza con un soporte finito orientado y con los lectores que conservan la
historia de sus rutas.  Su función es responder a una pregunta anterior a
cualquier comparación metrológica: ¿qué clases de persistencia puede sostener
la geometría nonádica y qué información debe conservarse para distinguirlas?

La respuesta tiene tres cardinales, pero no tres inventarios rivales:
\[
 81\ \text{celdas},\qquad
 56\ \text{familias de ruta},\qquad
 13\ \text{multisecciones familia--nivel}.
\]
El primer número cuenta posiciones orientadas; el segundo identifica las
posiciones que publican la misma firma finita de ruta; el tercero reúne las
posiciones que comparten familia y profundidad respecto del centro.  Ninguno
de ellos cuenta masas conocidas, nombres de especies o filas de una
comparación externa.

\section{La geometría completa de las \(81\) celdas}

Sea
\[
 I_9=\{1,\ldots,9\},
 \qquad
 \mathcal C_{81}=I_9\times I_9.
\]
Las coordenadas \((r,c)\) conservan las dos direcciones de APP\@. La paridad
de cada coordenada define cuatro sectores internos:
\[
 \operatorname{fam}(r,c)=
 \begin{cases}
  \ell^- ,& r\equiv1,\ c\equiv1\pmod2,\\
  \nu ,& r\equiv1,\ c\equiv0\pmod2,\\
  d ,& r\equiv0,\ c\equiv1\pmod2,\\
  u ,& r\equiv0,\ c\equiv0\pmod2.
 \end{cases}
 \tag{A1}\label{eq:familias-paridad-atlas}
\]
La profundidad en el atlas no se añade como etiqueta independiente.  Es la
distancia de Chebyshev al centro:
\[
 G(r,c)=\max\{|r-5|,|c-5|\}.
 \tag{A2}\label{eq:generacion-chebyshev-atlas}
\]
Por tanto, familia y nivel proceden de la propia posición.  En particular,
el centro no se obtiene seleccionando una especie desde fuera: es la única
celda de profundidad cero.

\begin{theorem}[Censo de familias y generaciones del atlas]
\label{thm:censo-atlas-81}
Las aplicaciones \eqref{eq:familias-paridad-atlas} y
\eqref{eq:generacion-chebyshev-atlas} descomponen el atlas de manera
disjunta como
\[
 81=45_{\rm leptones}+36_{\rm quarks}
   =25_{\ell^-}+20_{\nu}+20_{d}+16_{u}.
 \tag{A3}\label{eq:censo-81-cuatro-familias}
\]
Al refinar por \(G\), se obtienen exactamente las trece multisecciones
\[
\begin{array}{c|c|c}
\text{familia}&\text{niveles}&\text{cardinales}\\ \hline
\ell^-&G0,G2,G4&1,8,16\\
\nu&G1,G2,G3,G4&2,4,6,8\\
d&G1,G2,G3,G4&2,4,6,8\\
u&G1,G3&4,12.
\end{array}
\tag{A4}\label{eq:trece-multisecciones-atlas}
\]
La celda \((5,5)\) es la única componente de
\(\mathfrak S_{\ell^-,0}\).
\end{theorem}

\begin{proof}
En \(I_9\) hay cinco coordenadas impares y cuatro pares.  Las cuatro clases
de paridad tienen, por ello, cardinales
\[
 5\cdot5=25,\qquad
 5\cdot4=20,\qquad
 4\cdot5=20,\qquad
 4\cdot4=16,
\]
lo que prueba \eqref{eq:censo-81-cuatro-familias}; las dos primeras
clases reúnen \(45\) celdas y las dos últimas, \(36\).

Para el refinamiento, hacemos crecer cuadrados centrados en \((5,5)\).
En la clase impar--impar hay una celda en radio cero, \(3^2-1=8\)
en el cascarón de radio dos y \(5^2-3^2=16\) en el de radio cuatro.
En la clase impar--par, los rectángulos acumulados hasta radios
\(1,2,3,4\) contienen, respectivamente,
\[
 1\cdot2=2,\qquad 3\cdot2=6,\qquad
 3\cdot4=12,\qquad5\cdot4=20
\]
celdas; sus diferencias sucesivas son \(2,4,6,8\).  La clase
par--impar da el mismo cálculo con las coordenadas intercambiadas.  En la
clase par--par aparecen \(2^2=4\) celdas en radio uno y
\(4^2-2^2=12\) en radio tres.  Estos cascarones son disjuntos y agotan
\(I_9^2\), de modo que no queda otra multisección.
\end{proof}

La palabra \emph{nivel} designa en este capítulo la profundidad combinatoria
\(G\).  La lectura física posterior conserva esa geometría, pero no la
define.  Así se evita convertir los nombres de las partículas en premisas
del censo que debe explicarlas.

\section{Inversión central, \(41\) órbitas y unicidad del centro}

La inversión del atlas es
\[
 \rho_{\rm c}(r,c)=(10-r,10-c).
 \tag{A5}\label{eq:inversion-central-atlas-rev6}
\]
Como \(10-r\) tiene la misma paridad que \(r\), la inversión conserva las
cuatro familias.  Además,
\[
 |(10-r)-5|=|r-5|,
 \qquad
 |(10-c)-5|=|c-5|,
\]
y, por tanto, conserva el nivel \(G\).  Cada una de las trece
multisecciones es \(\rho_{\rm c}\)-estable.

\begin{proposition}[Órbitas y excepción central]
\label{prop:orbitas-centro-atlas-rev6}
La acción de \(\rho_{\rm c}\) sobre \(\mathcal C_{81}\) tiene exactamente
\(41\) órbitas.  Su único punto fijo es \((5,5)\).  En consecuencia, la
multisección \(\mathfrak S_{\ell^-,0}\) admite el único selector puntual
\(\rho_{\rm c}\)-equivariante; en las otras doce, el objeto natural es la
órbita o la multisección completa.
\end{proposition}

\begin{proof}
La ecuación
\[
 (r,c)=\rho_{\rm c}(r,c)
\]
equivale a \(2r=2c=10\), luego \(r=c=5\).  Las otras ochenta celdas se
agrupan en cuarenta pares, de donde \(1+40=41\) órbitas.

Sea \(y\) una multisección no central y supongamos que un selector puntual
equivariante elige \(a(y)\in\mathfrak S_y\).  Como \(\rho_{\rm c}\) actúa
trivialmente sobre el tipo familia--nivel, la equivarianza impondría
\[
 a(y)=a(\rho_{\rm c}y)=\rho_{\rm c}a(y).
\]
El punto elegido tendría que ser fijo, pero el único punto fijo está en
\(\mathfrak S_{\ell^-,0}\).  La contradicción demuestra la afirmación.
\end{proof}

Esta proposición explica por qué una especie interna no es, en general, una
casilla aislada.  El electrón central es excepcional porque su multisección
se reduce a un punto.  En las demás familias, la orientación conjugada forma
parte de la identidad que debe conservarse.

\section{Las \(56\) familias de ruta}

Cada celda porta la sombra finita de su historia de transporte:
\[
 \mathcal R(r,c)=
 \bigl(
 n_A,\ s_C,\ k_{\Delta_4},\
 \nu_{120},\ \nu_{270},\ \tau_{12}
 \bigr)_{r,c}.
 \tag{A6}\label{eq:lector-familia-ruta-atlas}
\]
Las cinco primeras coordenadas son enteras; \(\tau_{12}\) es la palabra
dodecafásica orientada.  Dos celdas pertenecen a la misma familia de ruta
si y sólo si tienen la misma sextupla \(\mathcal R\).

\begin{theorem}[Censo de familias de ruta y memoria mínima]
\label{thm:censo-56-familias-ruta-rev6}
La imagen de \(\mathcal R\) contiene exactamente \(56\) elementos.  El
histograma de tamaños de sus fibras es
\[
 41\times1,\qquad
 10\times2,\qquad
 2\times3,\qquad
 1\times4,\qquad
 2\times5.
 \tag{A7}\label{eq:histograma-fibras-ruta}
\]
En particular,
\[
 41+10+2+1+2=56,
\]
\[
 41\cdot1+10\cdot2+2\cdot3+1\cdot4+2\cdot5=81.
 \tag{A8}\label{eq:doble-censo-rutas}
\]
Fijado el orden lexicográfico dentro de cada fibra, el rango
\[
 \kappa(r,c)\in\{0,1,2,3,4\}
\]
hace inyectiva la aplicación
\[
 (r,c)\longmapsto\bigl(\mathcal R(r,c),\kappa(r,c)\bigr).
 \tag{A9}\label{eq:lift-memoria-cinco-atlas}
\]
Cinco es el tamaño mínimo de un alfabeto uniforme capaz de efectuar este
levantamiento.
\end{theorem}

\begin{proof}
La enumeración exhaustiva de \(I_9^2\) agrupa las ochenta y una sextuplas
\eqref{eq:lector-familia-ruta-atlas} por igualdad coordenada a coordenada.
Produce los cinco tamaños de fibra de
\eqref{eq:histograma-fibras-ruta}; las dos identidades
\eqref{eq:doble-censo-rutas} comprueban simultáneamente el número de clases
y la cobertura de las ochenta y una celdas.

El rango lexicográfico distingue los elementos de cada fibra, de modo que
\eqref{eq:lift-memoria-cinco-atlas} es inyectiva.  Como existen dos fibras
de cardinal cinco, cualquier alfabeto común con menos de cinco símbolos
identificaría dos celdas de una de ellas por el principio del palomar.
\end{proof}

El rango \(\kappa\) es memoria de posición dentro de una fibra; no es
carga, masa ni color.  Su función es recuperar la celda que la firma visible
de ruta había identificado con otras.  El número \(56\) cuenta clases de la
sextupla \(\mathcal R\), mientras el número \(13\) cuenta clases de la pareja
\((\operatorname{fam},G)\).  Son cocientes diferentes del mismo atlas.

\paragraph{Cocientes celular y orbital del atlas: dominios y mapas de proyección.}
La proyección cruda CT--108 y la clasificación científica leen datos
distintos de las mismas ochenta y una celdas. Si
\(q_{\rm CT}\) conserva solamente la forma de la palabra trítica
dodecafásica, mientras \(\mathcal R\) conserva la sextupla de ruta y
\(q=(\operatorname{fam},G)\) conserva familia y profundidad, el cuadro
tipado es
\[
\begin{array}{rcll}
 \mathcal C_{81}
 &\xrightarrow{\ q_{\rm CT}\ }&
 \mathcal W^{\rm raw}_{7}
 & \text{formas crudas CT--108},\\[2mm]
 \mathcal C_{81}
 &\xrightarrow{\ \mathcal R\ }&
 \Sigma^{\rm ruta}_{56}
 & \text{familias de ruta},\\[2mm]
 \mathcal C_{81}
 &\xrightarrow{\ q\ }&
 \Sigma_{13}
 & \text{multisecciones familia--nivel}.
\end{array}
\tag{A9a}\label{eq:dos-cocientes-atlas-rev8}
\]
Las siete fibras crudas tienen multiplicidades
\[
 54,2,2,2,7,7,7,
 \qquad54+3\cdot2+3\cdot7=81.
\]
Por tanto, el censo arqueológico \(81\to7\) no sustituye el atlas
\(81/56/13\). Tampoco la escritura abreviada \(81/56/13\) afirma que exista
una flecha canónica \(56\to13\): hay fibras de \(\mathcal R\) que atraviesan
más de una pareja familia--nivel. Los mapas \(\mathcal R\) y \(q\) son
cocientes paralelos y conservan invariantes diferentes. En el sector quark,
la convención activa sigue siendo
\[
 \operatorname{col}(r,c)=r-c\pmod3;
\]
la escritura histórica \(c-r\) es su convención conjugada, no otro color ni
otro cociente del atlas.

\section{Distinción entre celda, familia de ruta y multisección}

Definamos
\[
 q:\mathcal C_{81}\longrightarrow\Sigma_{13},
 \qquad
 q(r,c)=\bigl(\operatorname{fam}(r,c),G(r,c)\bigr),
 \tag{A10}\label{eq:cociente-trece-multisecciones}
\]
y, para \(y\in\Sigma_{13}\),
\[
 \mathfrak S_y^{\rm atlas}
 =q^{-1}(y)
 =\{(r,c)\in\mathcal C_{81}:q(r,c)=y\}.
 \tag{A11}\label{eq:multiseccion-atlas-rev6}
\]
Los cuatro niveles de lectura quedan entonces separados:

\begin{enumerate}[label=(\roman*)]
 \item una \emph{celda} es una posición orientada \((r,c)\), con todos sus
 datos locales;
 \item una \emph{familia de ruta} es una fibra de \(\mathcal R\): conserva
 una misma firma finita de transporte;
 \item una \emph{multisección} es una fibra de \(q\): conserva familia,
 profundidad y la totalidad de sus orientaciones conjugadas;
 \item una \emph{representación física} es la realización posterior de una
 multisección enriquecida en un espacio de estados; no es otra partición del
 tablero ni la elección arbitraria de una de sus celdas.
\end{enumerate}

Esta distinción da contenido matemático a la palabra \emph{especie}.  El tipo
interno es la multisección.  Una partícula concreta es una sección
persistente realizada en ella.  La representación física conserva la acción,
la orientación y los observables del tipo; no retrocede para redefinir el
atlas a partir de una masa ya conocida.

\section{Color residual como orientación ternaria}

En el sector quark, formado por las \(36\) celdas con \(r\) par, definimos
\[
 \operatorname{col}(r,c)=r-c\pmod3,
 \tag{A12}\label{eq:color-residual-atlas-rev6}
\]
con el diccionario
\[
 0\longmapsto R,\qquad1\longmapsto G,\qquad2\longmapsto B.
\]

\begin{proposition}[Balance de color residual]
\label{prop:color-atlas}
Las treinta y seis celdas de quark se reparten exactamente como
\[
 36=12_R+12_G+12_B.
 \tag{A13}\label{eq:balance-color-atlas-rev6}
\]
La inversión central fija \(R\) e intercambia \(G\) y \(B\).
\end{proposition}

\begin{proof}
Fijado cualquiera de los cuatro valores pares de \(r\), al recorrer
\(c=1,\ldots,9\) cada clase módulo tres aparece tres veces.  Cada residuo
tiene, por tanto, \(4\cdot3=12\) representantes.  Además,
\[
 \operatorname{col}\bigl(\rho_{\rm c}(r,c)\bigr)
 =(10-r)-(10-c)
 =-(r-c)\pmod3.
\]
La negación modular fija la clase cero e intercambia las clases uno y dos.
\end{proof}

El trítico \(R,G,B\) es, por tanto, un observable exacto del atlas y no una
decoración añadida a una tabla de quarks.  La representación de color se
construye después sobre este soporte ternario; el censo \(12+12+12\) es su
geometría previa.

\section{De la extensión superviviente a la firma}

El atlas clasifica cortes finitos.  La identidad de una partícula requiere la
historia profunda de la que ese corte procede.  Sea
\[
 \Omega_{\rm HMT}^{\rm surv}
 =\varprojlim_m\mathcal S_m
\]
el sistema inverso de estados supervivientes.  La proyección al atlas \(\pi_{81}\) y la sombra
de ruta forman la cadena
\[
 \Omega_{\rm HMT}^{\rm surv}
 \xrightarrow{\ \operatorname{sh}_{108}\ }
 \Gamma_{108}^{\rm enr}
 \xrightarrow{\ \mathcal E\ }
 \mathcal L_{108}
 \xrightarrow{\ L_{\rm tick}\ }
 \mathbb Z^5.
 \tag{A14}\label{eq:cadena-extension-firma-atlas-rev6}
\]
Aquí \(\operatorname{sh}_{108}\) publica una ruta de \(108\) pasos sin
olvidar hoja, orientación, frontera ni predecesores; \(\mathcal E\) forma el
registro dodecafásico; y \(L_{\rm tick}\) extrae
\[
 v(x)=
 \bigl(n_A,s_C,k_{\Delta_4},\nu_{120},\nu_{270}\bigr).
 \tag{A15}\label{eq:firma-cinco-extension-rev6}
\]
Ninguna de estas tres flechas consulta una masa, una unidad o el nombre de
una partícula.

La razón para conservar el registro antes de proyectar se ve en sus doce
eventos orientados \(e_0,\ldots,e_{11}\).  Para \(0\leq i\leq5\), sea
\[
 p_i=e_i+e_{i+6},
 \qquad
 a_i=e_i-e_{i+6},
 \qquad
 s_C=\sum_{i=0}^{5}p_i,
\]
y, para \(0\leq i<5\),
\[
 M_i=p_i-p_5.
\]
Entonces
\[
 p_5=\frac{s_C-\sum_{i=0}^{4}M_i}{6},
 \qquad
 p_i=p_5+M_i,
\]
\[
 e_i=\frac{p_i+a_i}{2},
 \qquad
 e_{i+6}=\frac{p_i-a_i}{2}.
 \tag{A16}\label{eq:reconstruccion-registro-doce}
\]
La identidad
\[
 1+5+6=12
\]
es así una reconstrucción reversible: un total, cinco diferencias y seis
orientaciones recuperan los doce eventos.  Proyectar a la firma de cinco
enteros antes de componer historias eliminaría precisamente la colocación y
la orientación que permiten decidir si dos rutas pueden componerse.

El cociente \(q\) se eleva a la historia profunda mediante la proyección al atlas
\(\pi_{81}:\Omega_{\rm HMT}^{\rm surv}\to\mathcal C_{81}\):
\[
 \widetilde{\mathfrak S}_y
 =\{x\in\Omega_{\rm HMT}^{\rm surv}:q(\pi_{81}x)=y\}.
 \tag{A17}\label{eq:multiseccion-enriquecida-rev6}
\]
Esta es la multisección enriquecida.  Su sombra pertenece al atlas; su ruta
pertenece a \(\Gamma_{108}^{\rm enr}\); su memoria completa pertenece al
registro; y su evaluación se aplica sólo al final.  La genealogía no es una
nota añadida al valor: es el objeto que hace posible el valor y conserva las
otras lecturas de la misma sección.

\section{Bifurcación del sector nulo antes del carácter positivo}

El carácter de masa actúa sobre la firma entera y tiene codominio positivo:
\[
 \chi_{\rm mass}:\mathbb Z^5\longrightarrow\mathbb R_{>0}.
 \tag{A18}\label{eq:caracter-positivo-atlas-rev6}
\]
Por definición, no existe una firma \(v\) con
\(\chi_{\rm mass}(v)=0\).  El sector nulo no es, por consiguiente, una
decimocuarta multisección ni el límite degenerado de las trece anteriores.

Su soporte algebraico pertenece a la rama partida
\[
 \mathcal A_{\rm sp}
 =\mathbb R[\mathsf R]/(\mathsf R^2-1),
 \qquad
 P_\pm=\frac{1\pm\mathsf R}{2}.
\]
Todo elemento se escribe de manera única como
\[
 z=r_+P_++r_-P_-,
 \qquad
 N_{\rm sp}(z)=r_+r_-,
\]
y las dos rectas
\[
 \mathbb R P_+,\qquad\mathbb R P_-
\]
son nulas porque \(N_{\rm sp}(\lambda P_\pm)=0\).  Una sección nula es una
ruta abierta transportada en una de esas direcciones.  El interior material
requiere el acoplamiento de ambas, para el cual \(r_+r_->0\).

Así quedan tipadas dos ramas:
\[
 \text{persistencia material}
 \longrightarrow\mathbb Z^5
 \xrightarrow{\chi_{\rm mass}}\mathbb R_{>0},
\]
\[
 \text{propagación nula}
 \longrightarrow\partial_{\rm null}\mathcal C_{\rm sp}^{+}.
\]
Las realizaciones fotónica y gauge leen después la segunda rama.  Su
separación del carácter positivo impide borrar fotón y gluón por no aparecer
en una lista de masas, y evita a la vez fabricar una masa nula mediante un
carácter cuyo codominio la excluye.

\section{Sistema finito de clasificación y memoria de trayectorias}

Los tres cocientes del capítulo cumplen funciones complementarias:
\[
 \mathcal C_{81}
 \xrightarrow{\ \mathcal R\ }\mathcal R_{56},
 \qquad
 \mathcal C_{81}
 \xrightarrow{\ q\ }\Sigma_{13},
 \qquad
 \mathcal C_{81}
 \xrightarrow{\ /\langle\rho_{\rm c}\rangle\ }
 \mathcal O_{41}.
\]
El primero conserva la firma de ruta; el segundo conserva familia y nivel;
el tercero conserva conjugación central.  Junto con el color residual y el
levantamiento \(\kappa\), forman un sistema finito capaz de reconstruir
qué información se ve, cuál se identifica y cuál debe permanecer como
memoria.

La cadena causal del atlas queda, por tanto, fijada:
\[
 \boxed{
 \begin{gathered}
 \text{extensión superviviente}
 \longrightarrow\text{ruta enriquecida}
 \longrightarrow\text{registro}
 \longrightarrow\text{firma}\\
 \longrightarrow\text{carácter}
 \longrightarrow\text{lectura dimensional}
 \end{gathered}}
 \tag{A19}\label{eq:cadena-rectora-atlas-rev6}
\]
Las ochenta y una celdas constituyen el dominio local; las cincuenta y seis
familias expresan coincidencias de ruta; las trece multisecciones son los
tipos internos; las cuarenta y una órbitas conservan la conjugación; y el
centro \((5,5)\) es la única excepción puntual.  La masa aparece después
como lectura de una persistencia ya construida.  El atlas no es una tabla
abreviada de resultados: es la estructura que explica por qué puede existir
una tabla.
